%! Polynomial Functions and Real Zeros — Unit 2, Lesson 2.3 — Group Activity (Answer Key)
\documentclass[10pt]{article}
\usepackage{precalculus-article}
\usepackage{precalculus-key}   % pulls in -boxes; do NOT also load -boxes
\newcommand{\TallMath}[1]{$\displaystyle #1\rule[-1.4em]{0pt}{3.2em}$}

\begin{document}

\pageheader{Unit 2, Lesson 2.3}{Group Activity --- Answer Key}

\vspace{0.12in}

\begin{scenariobox}[Four Polynomials, Two Zeros]{plumlight}
\small
Every polynomial below has exactly the same two real zeros, $x = -2$ and
$x = 1$. Only the \textbf{multiplicities} change --- and that is enough to
change the picture. All three graphs are pre-drawn; read them, do not draw.

\smallskip
\begin{center}
\begin{tabular}{@{}llll@{}}
\textbf{A:} $p(x) = (x+2)(x-1)$ &\quad
\textbf{B:} $p(x) = (x+2)(x-1)^2$ &\quad
\textbf{C:} $p(x) = (x+2)^2(x-1)^2$ &\quad
\textbf{D:} $p(x) = (x+2)^2(x-1)^3$
\end{tabular}
\end{center}

\smallskip
\begin{center}
\begin{tabular}{@{}ccc@{}}
\begin{tikzpicture}[x=0.52cm, y=0.4cm]
  \draw[linegray, very thin, step=1] (-3,-2) grid (3,3);
  \draw[->, thick] (-3.3,0) -- (3.3,0) node[below right] {\scriptsize $x$};
  \draw[->, thick] (0,-2.3) -- (0,3.3);
  \foreach \x in {-2,-1,1,2} \node[below] at (\x,-0.15) {\scriptsize $\x$};
  \draw[very thick, plum] plot [smooth] coordinates
    {(-2.6,-1.94) (-2.2,-0.51) (-2,0) (-1.5,0.78) (-1,1) (-0.5,0.84)
     (0,0.5) (0.5,0.16) (1,0) (1.5,0.22) (2,1) (2.4,2.16)};
  \fill[plum] (-2,0) circle (2pt);
  \fill[plum] (1,0) circle (2pt);
\end{tikzpicture}
&
\begin{tikzpicture}[x=0.52cm, y=0.4cm]
  \draw[linegray, very thin, step=1] (-3,-2) grid (3,3);
  \draw[->, thick] (-3.3,0) -- (3.3,0) node[below right] {\scriptsize $x$};
  \draw[->, thick] (0,-2.3) -- (0,3.3);
  \foreach \x in {-2,-1,1,2} \node[below] at (\x,-0.15) {\scriptsize $\x$};
  \draw[very thick, plum] plot [smooth] coordinates
    {(-3,2) (-2.5,0.875) (-2,0) (-1.5,-0.625) (-1,-1) (-0.5,-1.125)
     (0,-1) (0.5,-0.625) (1,0) (1.5,0.875) (2,2)};
  \fill[plum] (-2,0) circle (2pt);
  \fill[plum] (1,0) circle (2pt);
\end{tikzpicture}
&
\begin{tikzpicture}[x=0.52cm, y=0.4cm]
  \draw[linegray, very thin, step=1] (-3,-2) grid (3,3);
  \draw[->, thick] (-3.3,0) -- (3.3,0) node[below right] {\scriptsize $x$};
  \draw[->, thick] (0,-2.3) -- (0,3.3);
  \foreach \x in {-2,-1,1,2} \node[below] at (\x,-0.15) {\scriptsize $\x$};
  \draw[very thick, plum] plot [smooth] coordinates
    {(-2.7,1.68) (-2.5,0.77) (-2,0) (-1.5,0.39) (-1,1) (-0.5,1.27)
     (0,1) (0.5,0.39) (1,0) (1.5,0.77) (1.7,1.68)};
  \fill[plum] (-2,0) circle (2pt);
  \fill[plum] (1,0) circle (2pt);
\end{tikzpicture}
\\[2pt]
\textbf{Graph I} & \textbf{Graph II} & \textbf{Graph III}
\end{tabular}
\end{center}
\end{scenariobox}

\vspace{0.1in}

\boxguard[30]
\begin{tcolorbox}[colback=white, colframe=black!40,
    title=\textbf{Tier R --- Remediate},
    fonttitle=\bfseries, arc=2mm, left=3mm, right=3mm, top=2mm, bottom=2mm]

\begin{enumerate}[itemsep=8pt]
  \item Complete the table for polynomials A--D. The degree is the sum of the
        multiplicities.

        \smallskip
        \renewcommand{\arraystretch}{1.6}
        \begin{tabular}{|c|c|c|c|}
        \hline
        & \textbf{Multiplicity of $x = -2$} & \textbf{Multiplicity of $x = 1$} & \textbf{Degree} \\
        \hline
        \textbf{A} & \ans{1} & \ans{1} & \ans{2} \\
        \hline
        \textbf{B} & \ans{1} & \ans{2} & \ans{3} \\
        \hline
        \textbf{C} & \ans{2} & \ans{2} & \ans{4} \\
        \hline
        \textbf{D} & \ans{2} & \ans{3} & \ans{5} \\
        \hline
        \end{tabular}

  \item Circle the correct word. At $x = 1$, polynomial \textbf{A}
        \ans{crosses} / \textbf{is tangent to} the $x$-axis, and
        polynomial \textbf{B} \textbf{crosses} / \ans{is tangent to} the
        $x$-axis.

  \item Which polynomial is tangent to the $x$-axis at \textit{both}
        intercepts? \quad \ans{C}
\end{enumerate}
\end{tcolorbox}

\vspace{0.1in}

\boxguard
\begin{tcolorbox}[colback=white, colframe=black!40,
    title=\textbf{Tier A --- Approaching Proficiency},
    fonttitle=\bfseries, arc=2mm, left=3mm, right=3mm, top=2mm, bottom=2mm]

\begin{enumerate}[itemsep=8pt]
  \item Match each graph to one of A--D. One formula is left over.

        \smallskip
        Graph I: \ans{B} \qquad Graph II: \ans{A} \qquad
        Graph III: \ans{C} \qquad Left over: \ans{D}

  \item Justify \textbf{one} of your matches using multiplicity --- not by
        the shape alone. Name the intercept, its multiplicity, and what the
        graph does there.

        \vspace{0.05in}
        \ansline{Graph I is B: at $x = 1$ the factor $(x-1)$ is squared, an even multiplicity,}
        \ansline{so the graph touches and turns back; at $x = -2$ multiplicity 1 means it crosses.}

  \item The leftover formula is not graphed above. Fill in what its graph
        \textit{would} do:

        \smallskip
        At $x = -2$ the multiplicity is \ans{2}, so the graph
        \ans{is tangent to} the axis. At $x = 1$ the multiplicity is
        \ans{3}, so the graph \ans{crosses} the axis.

  \item Exactly one of A--D has \textit{both} multiplicities even, so its
        graph touches the $x$-axis twice and never passes through it. Which
        one is it, and how do the multiplicities tell you?

        \vspace{0.05in}
        \ansline{C: both $(x+2)$ and $(x-1)$ appear twice, so both multiplicities are even.}
        \ansline{Even multiplicity means the outputs on both sides share a sign --- no crossing.}
\end{enumerate}
\end{tcolorbox}

\vspace{0.1in}

\boxguard
\begin{tcolorbox}[colback=white, colframe=black!40,
    title=\textbf{Tier E --- Extension},
    fonttitle=\bfseries, arc=2mm, left=3mm, right=3mm, top=2mm, bottom=2mm]

\begin{enumerate}[itemsep=8pt]
  \item Build a polynomial to order: degree 3, a zero at $x = -3$ of
        multiplicity 2, a zero at $x = 1$, and $p(0) = 18$.

        \begin{work}
          p(x) &= a(x+3)^2(x-1) \\
          p(0) &= a(3)^2(-1) \\
               &= -9a \\
           -9a &= 18 \\
             a &= -2
        \end{work}

        $p(x) = $ \ans{$-2(x+3)^2(x-1)$}

  \item Write a degree-4 polynomial whose graph meets the $x$-axis at exactly
        two points and \textit{never} crosses it. Then explain, in terms of
        multiplicity, why your answer works.

        \vspace{0.05in}
        $p(x) = $ \ans{$(x-1)^2(x+2)^2$}

        \vspace{0.05in}
        \ansline{Each zero has multiplicity 2 --- even --- so the graph is tangent at both}
        \ansline{$x = 1$ and $x = -2$; $2 + 2 = 4$ gives the required degree.}

  \item The table below shows a polynomial $q$ at equally spaced inputs.

        \smallskip
        \renewcommand{\arraystretch}{1.4}
        \begin{tabular}{|c|c|c|c|c|c|c|}
        \hline
        $x$ & 0 & 1 & 2 & 3 & 4 & 5 \\
        \hline
        $q(x)$ & 0 & 2 & 16 & 54 & 128 & 250 \\
        \hline
        \end{tabular}
        \smallskip

        \begin{enumerate}[label=(\alph*), itemsep=4pt]
          \item The first constant row of differences is the \ans{3rd}
                differences, so $\deg q = $ \ans{3}
          \item $q(0) = 0$, so $x = 0$ is a zero. Could that zero have
                multiplicity 4? Explain using the degree.

                \vspace{0.05in}
                \ansline{No. A factor $x^4$ alone would make the degree at least 4, but the third}
                \ansline{differences are constant, so $\deg q = 3$. Multiplicity can never exceed degree.}
        \end{enumerate}
\end{enumerate}
\end{tcolorbox}

\end{document}
